\chapter{Python à grande échelle}\label{ch:pl:python-at-scale}

Une première version du signal d'un chercheur parcourt les événements de la journée un par un~: pour chaque enregistrement, il lit son côté et sa quantité,
met à jour une moyenne mobile exponentielle, puis la stocke. Sur l'ordinateur portable, cela coûte 960~nanosecondes par événement, soit trois minutes pour une journée de
deux cents millions de messages. Le même calcul, écrit comme une seule expression sur tableau, coûte 14~nanosecondes par événement, et, traité
en flux par blocs de dix mille événements, il ne retient jamais plus d'un tiers de mégaoctet, quelle que soit la longueur de
l'historique. Les mathématiques n'ont pas changé. Ce qui a changé, c'est le travail que l'interpréteur effectue par nombre, le nombre d'octets
que coûte chaque représentation des données, et ce qui se trouve en mémoire à un instant donné. Ce chapitre mesure ces trois coûts sur des données de la
\omterm{def:pl:a-tick-store-on-open-formats:store}{base de ticks} et transpose ces enseignements dans l'outil de \omterm{def:pl:python-at-scale:chunked}{traitement par blocs} de la plateforme de recherche (\texttt{firm.pyscale}).

\section{Ce que coûte l'interpréteur}

Python est le langage de la plateforme de recherche parce qu'il est rapide à écrire et qu'il dispose des bibliothèques~; le prix à payer est que chaque
opération écrite en Python est exécutée par un interpréteur qui manipule des objets.

\begin{definition}[Surcoût de l'interpréteur]\label{def:pl:python-at-scale:overhead}
Le \emph{surcoût de l'interpréteur}\index{surcoût de l'interpréteur} d'une opération est le travail que l'interpréteur Python effectue autour de l'arithmétique
elle-même~: décodage du bytecode, recherche des noms et des attributs, vérification des types, création et destruction des objets
qui portent les valeurs intermédiaires, et comptage des références. Pour une opération scalaire, il vaut de dix à cent fois le coût
de l'arithmétique.
\end{definition}

Le \Cref{lst:pl:python-at-scale:loops} est la première version du signal, en deux exemplaires. La première boucle lit les enregistrements projetés en mémoire de
la \omterm{def:pl:a-tick-store-on-open-formats:store}{base de ticks}, un par un~; chaque accès à un champ de \texttt{rec[i]} crée un objet enregistrement, puis un objet scalaire. La seconde convertit
une fois pour toutes les deux colonnes en listes Python, puis boucle sur des entiers ordinaires.

\omcode{code/platforms/08-python-at-scale/python/pl_pyscale.py}{53}{73}{Le signal sous forme de boucle Python~: sur les enregistrements projetés, et
sur des listes construites une fois à partir des colonnes.}\label{lst:pl:python-at-scale:loops}

La \Cref{fig:pl:python-at-scale:methods} les mesure~: 960~nanosecondes par événement sur les enregistrements, 145 sur les listes. L'écart
est le coût de la construction des objets que crée l'accès aux enregistrements~; ce qui reste, 150~nanosecondes pour une dizaine d'opérations de bytecode,
c'est l'interpréteur lui-même. Aucun des deux n'est un bogue~; ce sont tous deux le prix du langage.

\begin{omfigure}
\begin{tikzpicture}
\begin{axis}[width=11cm,height=5cm,xbar,bar width=9pt,xmode=log,log origin=infty,xmin=5,xmax=3000,
  xlabel={nanosecondes par événement (échelle logarithmique)},y dir=reverse,
  ytick={0,1,2,3},yticklabels={boucle sur les enregistrements,boucle sur les listes,expression de tableau,par blocs~; 65\,536},
  tick label style={font=\small},label style={font=\small},grid=major,grid style={gray!20},point meta=explicit,
  nodes near coords,nodes near coords style={font=\footnotesize,anchor=west,/pgf/number format/fixed,/pgf/number format/precision=1},
  table/col sep=comma]
\addplot[fill=omDef!45,draw=omDef] table[x=ns_per_event,y expr=\coordindex,meta=ns_per_event]{figdata/platforms/08-python-at-scale/measured_methods.csv};
\end{axis}
\end{tikzpicture}

\omcaption{Le coût par événement d'un même signal (la moyenne mobile exponentielle du volume signé) calculé de quatre façons sur des
enregistrements de la \omterm{def:pl:a-tick-store-on-open-formats:store}{base de ticks}~: une boucle Python sur les enregistrements projetés, une boucle Python sur des listes, une expression
sur tableau, et la même expression traitée en flux par blocs de 65\,536 événements. Mesuré sur un ordinateur portable (Intel Core Ultra 7 155H) sous WSL2, un seul fil d'exécution,
machine par ailleurs inactive. Données~: \texttt{bench\_pyscale.py}.}\label{fig:pl:python-at-scale:methods}
\end{omfigure}

\section{Penser en tableaux}

\begin{definition}[Programmation par tableaux]\label{def:pl:python-at-scale:array}
La \emph{programmation par tableaux}\index{programmation par tableaux} exprime un calcul comme des opérations sur des tableaux entiers --- arithmétique
élément par élément, réductions, sommes cumulées, filtres, tris --- de sorte que l'interpréteur n'est sollicité qu'une fois par opération et que la boucle
sur les éléments s'exécute en code compilé sur une mémoire contiguë.
\end{definition}

La seule difficulté réelle est la récurrence. Une moyenne mobile exponentielle $y_t = (1-\alpha)\,y_{t-1} + \alpha\,x_t$ semble intrinsèquement
séquentielle, mais c'est un filtre linéaire~: sa sortie est l'entrée convoluée avec $\alpha(1-\alpha)^k$, et une routine de filtrage la calcule
en code compilé en un seul appel (le \cref{lst:pl:python-at-scale:ewma}). Sur l'ordinateur portable, cela fait 14~nanosecondes par événement, soit 70 fois
plus vite que la boucle sur les enregistrements et 11 fois plus vite que la boucle sur les listes, et les résultats sont identiques jusqu'au dernier bit~:
le filtre effectue les mêmes multiplications et additions, dans le même ordre.

\omcode{code/firm/pyscale/firm_pyscale.py}{39}{50}{La moyenne récurrente sous forme de filtre linéaire~; l'état interne du filtre pour le
premier élément provient de la dernière sortie du bloc précédent.}\label{lst:pl:python-at-scale:ewma}

\begin{method}[De la boucle aux tableaux]\label{met:pl:python-at-scale:arrays}
Repérer le motif dans la boucle~: une application élément par élément (arithmétique, \texttt{where})~; un total cumulé (\texttt{cumsum})~; une fenêtre
(différences de sommes cumulées, ou vue glissante)~; une récurrence linéaire (un filtre)~; un regroupement (tri, puis réduction par segments)~;
une recherche «~as-of~» (\texttt{searchsorted}, chapitre~4). Ce qui ne s'y prête pas --- un carnet d'ordres qui change à chaque événement, une machine
à états aux nombreux embranchements --- relève du code compilé (chapitre~9), et non d'une boucle plus lente.
\end{method}

Un compilateur qui transforme la boucle Python en code machine à l'exécution est l'autre voie, la compilation à la volée du
livre~13 (chapitre~10)~; l'environnement de la série n'en installe pas, et le chapitre~9 emprunte la voie explicite d'écrire le noyau
en C++ ou en Rust.

\section{Mémoire~: objets, tableaux et copies}

Le second coût est l'espace. Un objet Python porte un en-tête (un pointeur de type et un compteur de références) avant sa valeur, et une liste de
tels objets est un tableau de pointeurs vers des objets dispersés en mémoire. Un tableau de nombres machine, lui, ne contient que les nombres.

\begin{omfigure}
\begin{tikzpicture}
\begin{axis}[width=11cm,height=4.8cm,xbar,bar width=9pt,xmin=0,xmax=72,xlabel={mégaoctets pour un million de symboles},y dir=reverse,
  ytick={0,1,2,3},yticklabels={object,category,arrow string,numpy à largeur fixe},
  tick label style={font=\small},label style={font=\small},grid=major,grid style={gray!20},
  nodes near coords,nodes near coords style={font=\footnotesize,/pgf/number format/fixed,/pgf/number format/precision=1},
  table/col sep=comma]
\addplot[fill=omMeth!45,draw=omMeth] table[x=mb,y=k]{figdata/platforms/08-python-at-scale/memory.csv};
\end{axis}
\end{tikzpicture}

\omcaption{La même colonne d'un million de symboles (cinquante codes boursiers distincts de quatre caractères) sous quatre représentations~: objets pandas
(une chaîne Python par ligne), catégorielle pandas (codes entiers et cinquante chaînes), colonne de chaînes Arrow (décalages et octets), et
tableau numpy à largeur fixe en Unicode. Données~: \texttt{fig\_pyscale.py}.}\label{fig:pl:python-at-scale:memory}
\end{omfigure}

La \Cref{fig:pl:python-at-scale:memory} mesure une colonne de symboles. Sous forme d'objets Python, elle occupe 61~Mo, soit un objet chaîne par ligne, dont
quelques octets sont les caractères~; sous forme de catégorielle, 1.0~Mo, un petit entier par ligne et cinquante chaînes une seule fois~; en chaînes Arrow,
8~Mo, un décalage et quatre octets par ligne~; sous forme de tableau numpy à largeur fixe en Unicode, 16~Mo, quatre caractères de quatre octets chacun.
La catégorielle est 61 fois plus petite que les objets~: l'\omterm{def:pl:columnar-formats:encodings}{encodage par dictionnaire} du chapitre~3, en mémoire. Toute colonne
d'un dataframe de recherche qui contient des chaînes répétées (symboles, plateformes, conditions, côtés) mérite la conversion.

\begin{remark}[Copies]\label{rem:pl:python-at-scale:copies}
La \omterm{def:pl:python-at-scale:array}{programmation par tableaux} a son propre piège~: chaque expression intermédiaire alloue un nouveau tableau.
\texttt{np.where(side == B, 1.0, -1.0) * qty} sur une journée de deux cents millions d'événements alloue un tableau de booléens, deux tableaux de flottants et un produit, soit plusieurs gigaoctets de
temporaires pour un résultat de 1.6~Go. Les remèdes consistent à calculer sur place là où la bibliothèque le permet, à choisir des types étroits et,
surtout, à ne pas garder la journée entière en mémoire à la fois.
\end{remark}

\section{Blocs et boucles hors mémoire}

\begin{definition}[Traitement par blocs, pic de mémoire résidente]\label{def:pl:python-at-scale:chunked}
Le \emph{traitement par blocs}\index{traitement par blocs} exécute un calcul sur un jeu de données tranche par tranche, en ne transportant d'une tranche à l'autre
que l'état nécessaire pour poursuivre le calcul à l'identique (la dernière valeur d'un filtre, la queue d'une fenêtre), de sorte que la mémoire utilisée
soit fixée par la taille du bloc et non par la taille des données. Le \emph{pic de mémoire résidente}\index{pic de mémoire résidente} d'un processus est
la plus grande quantité de mémoire physique qu'il a occupée à un instant donné~; c'est elle, et non la moyenne, qui décide si un travail tient sur une machine.
\end{definition}

Les fichiers plats du magasin de ticks sont projetés en mémoire (chapitre~4), si bien qu'une tranche du fichier est une vue~: rien n'est lu tant qu'elle n'est pas touchée,
et rien n'a besoin d'être libéré. Le \Cref{lst:pl:python-at-scale:stream} fait défiler le signal à travers le fichier et ne conserve que son
maximum courant~; l'état du filtre traverse les frontières entre blocs, de sorte que le résultat est égal à celui du calcul sur le tableau entier.

\omcode{code/platforms/08-python-at-scale/python/pl_pyscale.py}{152}{159}{Le signal diffusé à travers le fichier projeté en mémoire~: les temporaires d'un bloc
et l'état du filtre sont les seuls éléments jamais présents en mémoire.}\label{lst:pl:python-at-scale:stream}

\begin{proposition}[Bloc par bloc égal au tout]\label{prop:pl:python-at-scale:exact}
Si un calcul sur une suite peut s'écrire $y_t = f(y_{t-1}, x_t)$, ou comme une fonction d'une fenêtre des $w$ dernières entrées, alors
traiter la suite par blocs, en transmettant la dernière sortie (ou les $w-1$ dernières entrées) de chaque bloc au suivant, donne les mêmes
sorties que le traitement de la suite d'un seul tenant.
\end{proposition}

\begin{proof}
Par récurrence sur les blocs~: la première sortie d'un bloc est $f$ appliquée à l'état transmis et à sa première entrée,
ce que le calcul complet applique à cette position~; à l'intérieur d'un bloc, la récurrence est la même.
\end{proof}

\begin{omfigure}
\begin{tikzpicture}
\begin{axis}[width=11cm,height=5.6cm,ymode=log,xlabel={taille du bloc, événements},ylabel={pic de mémoire, Mo (échelle log)},
  tick label style={font=\small},label style={font=\small},grid=major,grid style={gray!20},xmin=-0.8,xmax=4.4,ymin=0.01,ymax=300,
  xtick={0,1,2,3,4},xticklabels={1\,000,10\,000,100\,000,1\,000\,000,entier},ytick={0.01,0.1,1,10,100},
  yticklabels={0.01,0.1,1,10,100},
  nodes near coords,point meta=explicit,nodes near coords style={font=\footnotesize,anchor=south east,
  /pgf/number format/fixed,/pgf/number format/precision=3},table/col sep=comma]
\addplot[thick,mark=*,omDef] table[x expr=\coordindex,y=peak_mb,meta=seconds]{figdata/platforms/08-python-at-scale/measured_chunks.csv};
\end{axis}
\end{tikzpicture}

\omcaption{Pic de mémoire du signal diffusé sur un historique de deux millions d'événements en fonction de la taille du bloc, chaque point étiqueté avec
sa durée en secondes~; le dernier point calcule le tableau entier d'un seul coup. La mémoire croît avec le bloc~; le temps est minimal pour des blocs de
cent mille événements, et il reste à moins de dix pour cent de ce minimum de dix mille événements jusqu'au tableau entier, une fois le surcoût
par bloc amorti. Mesuré sur un ordinateur portable
(Intel Core Ultra 7 155H) sous WSL2, un seul fil d'exécution, machine par ailleurs inactive~; mémoire comptée par tracemalloc, l'entrée projetée
non comprise. Données~: \texttt{bench\_pyscale.py}.}\label{fig:pl:python-at-scale:chunks}
\end{omfigure}

La \cref{fig:pl:python-at-scale:chunks} illustre le compromis. Des blocs de mille événements paient l'interpréteur une fois par mille et prennent
0.048~s pour les deux millions d'événements~; des blocs de dix mille prennent 0.025~s avec un pic de 0.31~Mo, et des blocs de cent mille,
les plus rapides ici, 0.023~s avec 2.5~Mo~; le tableau entier d'un seul coup prend 0.028~s avec un pic de 32~Mo, soit 16~octets par événement. À 16~octets par événement, la règle de la série, qui limite à 1.5~Go un processus, autorise environ
94 millions d'événements d'un seul coup, une petite fraction d'une journée d'un grand flux~; en mode diffusé, la seule limite est le temps.

\section{Processus, fils d'exécution et verrou de l'interpréteur}

\begin{definition}[Verrou global de l'interpréteur]\label{def:pl:python-at-scale:gil}
Le \emph{verrou global de l'interpréteur}\index{verrou global de l'interpréteur} (global interpreter lock) de l'interpréteur Python standard n'autorise qu'un seul fil d'exécution à la fois à
exécuter du bytecode Python~; les fils d'exécution exécutent du code Python de façon concurrente mais non en parallèle, tandis que le code compilé qui libère le verrou (une grande partie
de numpy, entrées-sorties) s'exécute en parallèle avec eux.
\end{definition}

Deux fils d'exécution faisant chacun une boucle Python pure de cinq millions d'additions prennent 0.20~s, contre 0.11~s pour un seul~: deux fois le travail en deux fois
le temps, aucun parallélisme. Deux fils d'exécution triant chacun un tableau numpy de huit millions de valeurs prennent 0.094~s contre 0.086~s pour un seul~: le
tri libère le verrou, et les deux s'exécutent en grande partie en parallèle. Le code Python qui doit utiliser plusieurs cœurs recourt donc à plusieurs processus.

\begin{dated}{2026-09}{Un interpréteur sans verrou}\label{dat:pl:python-at-scale:freethreading}
La PEP~703 (statut Final, consultée en septembre 2026) a ajouté une configuration de compilation de CPython qui exécute du code Python sans le
\omterm{def:pl:python-at-scale:gil}{verrou global de l'interpréteur}~; à partir de la version 3.13, CPython prend en charge cette version free-threaded, et les installeurs officiels pour macOS et Windows
peuvent l'installer en option. Le verrou reste l'option par défaut. L'environnement de la série (Python 3.10) possède le verrou.
\end{dated}

\begin{definition}[Parallélisme par processus, coût de sérialisation]\label{def:pl:python-at-scale:processes}
Le \emph{parallélisme par processus}\index{parallélisme par processus} exécute le travail dans plusieurs processus du système d'exploitation, chacun avec son propre
interpréteur et son propre verrou, de sorte que le code Python s'exécute sur plusieurs cœurs à la fois. Son \emph{coût de sérialisation}\index{coût de sérialisation} est
le temps passé à convertir les données en octets puis à les reconvertir pour les déplacer d'un processus à l'autre (sérialisation par pickle, copie par un tube), coût que la mémoire partagée
ou les fichiers projetés par les deux processus permettent d'éviter.
\end{definition}

Déplacer un tableau de 64~Mo vers un second processus et en sommer le contenu prend 0.14~s par pickle à travers un tube et 0.03~s par mémoire partagée,
la copie dans le bloc partagé comprise~: la sérialisation domine dans le cas du tube, et une somme de huit millions de nombres est négligeable
devant l'une ou l'autre. La règle de la plateforme de recherche en découle~: déplacer le travail vers les données, et non les données vers les processus de calcul. Confier à chaque processus une tranche
d'un \omterm{def:pl:a-tick-store-on-open-formats:mmap}{fichier projeté en mémoire}, une liste de partitions Parquet ou un bloc partagé, et lui faire renvoyer un petit résultat.

\begin{omfigure}
\begin{tikzpicture}
\begin{axis}[width=11cm,height=5cm,xbar,bar width=9pt,xmin=0,xmax=0.27,xlabel={secondes},y dir=reverse,
  xtick={0,0.05,0.10,0.15,0.20,0.25},xticklabels={0,0.05,0.10,0.15,0.20,0.25},
  ytick={0,1,2,3,4,5},yticklabels={sérialisé par un tube,mémoire partagée,boucle Python~; 1 thread,boucle Python~; 2 threads,tri numpy~; 1 thread,tri numpy~; 2 threads},
  tick label style={font=\small},label style={font=\small},grid=major,grid style={gray!20},
  nodes near coords,nodes near coords style={font=\footnotesize,/pgf/number format/fixed,/pgf/number format/precision=3},
  table/col sep=comma]
\addplot[fill=omThm!45,draw=omThm] table[x=seconds,y expr=\coordindex]{figdata/platforms/08-python-at-scale/measured_parallel.csv};
\end{axis}
\end{tikzpicture}

\omcaption{Déplacement d'un tableau de 64~Mo vers un second processus (sérialisé par pickle à travers un tube, ou placé en mémoire partagée), et exécution d'une boucle Python pure
et d'un tri numpy sur un fil d'exécution puis sur deux (chaque fil exécutant la tâche entière). Mesuré sur un ordinateur portable (Intel Core Ultra 7 155H)
sous WSL2, machine par ailleurs inactive. Données~: \texttt{bench\_pyscale.py}.}\label{fig:pl:python-at-scale:parallel}
\end{omfigure}

\section{Tutoriel~: quatre façons de calculer un signal}\label{tut:pl:python-at-scale:tut}

\begin{tutorial}
\textbf{Objectif.} Calculer un signal sur des enregistrements de la \omterm{def:pl:a-tick-store-on-open-formats:store}{base de ticks} par une boucle, par des tableaux et par blocs, vérifier que les réponses concordent, et
mesurer le temps et la mémoire. \textbf{État final~:} les \cref{fig:pl:python-at-scale:methods,fig:pl:python-at-scale:memory,fig:pl:python-at-scale:chunks}.
\begin{tutsteps}
\item \textbf{Historique}~: \texttt{pl\_pyscale.history(path, n)} écrit $n$ événements de version 1 sous forme de fichier plat de la \omterm{def:pl:a-tick-store-on-open-formats:store}{base de ticks}.
\item \textbf{Concordance}~: \texttt{agree()} exécute les deux boucles (\cref{lst:pl:python-at-scale:loops}), l'expression sur tableaux et la
version par blocs sur 20\,000 événements~: les écarts sont nuls.
\item \textbf{Mesure}~: \texttt{bench\_pyscale.py} chronomètre les quatre méthodes, les tailles de bloc et les tâches parallèles.
\item \textbf{Mémoire}~: \texttt{symbol\_memory()} pour les quatre représentations d'une colonne.
\item \textbf{Fenêtres}~: \texttt{firm\_pyscale.rolling\_sum\_kernel(w)} transporte la queue d'une fenêtre d'un bloc à l'autre~; il faut en vérifier le résultat contre une
somme cumulée calculée d'un seul tenant.
\end{tutsteps}
\textbf{Pour aller plus loin.} Diffuser le signal à partir des partitions Parquet de la semaine, décrites au chapitre~4, au lieu du fichier plat~; répartir le
l'historique entre deux processus par moitiés et combiner leurs résultats, et repérer quel état franchit la frontière.
\end{tutorial}

\section{Construction~: le pipeline par blocs}\label{bld:pl:python-at-scale:pyscale}

\begin{build}
\bfield{Rôle} La manière dont la plateforme de recherche exécute tout calcul par événement sur des historiques plus grands que la mémoire, et mesure le coût
en temps et en mémoire d'un calcul avant sa planification (chapitre~13).
\bfield{Interface} \texttt{chunks(array, size)}~; \texttt{run(source, kernel, state) -> (outputs, state)}, un noyau étant
\texttt{kernel(chunk, state) -> (output, state)}~; \texttt{ewma\_kernel(alpha)}, \texttt{rolling\_sum\_kernel(window)}~;
\texttt{probe(fn, *args) -> (result, seconds, peak\_bytes)}.
\bfield{Règles} Les blocs sont des vues, jamais des copies~; l'état transporté est le plus petit qui rend le résultat par blocs égal au
résultat complet, et un test le prouve pour chaque noyau~; la mémoire se mesure comme un pic, la méthode étant énoncée.
\bfield{Tests d'acceptation} \texttt{code/firm/pyscale/tests/}~: les blocs couvrent le tableau exactement~; les noyaux EWMA et de somme glissante donnent le
résultat complet pour toute taille de bloc, y compris un et la longueur totale~; la sonde signale un pic qui croît avec l'allocation.
\bfield{Pour aller plus loin} Un pool de processus de travail recevant chacun un intervalle d'un \omterm{def:pl:a-tick-store-on-open-formats:mmap}{fichier projeté en mémoire} (memory-mapped) (deux au plus ici, par la règle des ressources)~; un noyau
pour un filtre de signe de la transaction qui exige une machine à états~; le \omterm{def:pl:python-at-scale:chunked}{pic de mémoire résidente} fourni par le système d'exploitation à côté du décompte de tracemalloc.
\end{build}

\begin{omsources}
\item S.~Gross, \emph{PEP 703 -- Making the Global Interpreter Lock Optional in CPython} (Final)~; documentation Python,
\emph{Python support for free threading}.
\item Documentation Python, \texttt{multiprocessing.shared\_memory}.
\end{omsources}

\section{Exercices}

\begin{exercise}[$\star$]\label{exo:pl:python-at-scale:1}
À 960, 145 et 14 nanosecondes par événement, combien de temps chaque méthode met-elle pour une journée de deux cents millions d'événements~?
\end{exercise}

\begin{exercise}[$\star$]\label{exo:pl:python-at-scale:2}
Pourquoi la boucle sur les enregistrements coûte-t-elle environ six fois la boucle sur les listes, alors que les deux exécutent la même arithmétique~?
\end{exercise}

\begin{exercise}[$\star$]\label{exo:pl:python-at-scale:3}
Un dataframe compte un million de lignes et trois colonnes de chaînes de cinquante valeurs distinctes chacune. Estimer sa mémoire de chaînes en objets et
en catégories, à partir de la \cref{fig:pl:python-at-scale:memory}.
\end{exercise}

\begin{exercise}[$\star\star$]\label{exo:pl:python-at-scale:4}
Écrire la somme glissante des $w$ dernières valeurs comme une expression sur tableaux, et énoncer l'état qu'une version par blocs doit transporter.
\end{exercise}

\begin{exercise}[$\star\star$]\label{exo:pl:python-at-scale:5}
Deux threads ont exécuté la même boucle Python pure en deux fois le temps d'une seule. Qu'auraient fait deux processus, et que
coûteraient-ils si chacun avait besoin d'une entrée de 64~Mo~?
\end{exercise}

\begin{exercise}[$\star\star$]\label{exo:pl:python-at-scale:6}
À 16 octets par événement pour le calcul sur le tableau entier, combien d'événements tiennent d'un coup dans un budget de 1.5~Go, et quelle fraction
d'une journée du flux d'options prévu de 311 milliards de messages (chapitre~2) cela représente-t-il~?
\end{exercise}

\begin{exercise}[$\star\star\star$]\label{exo:pl:python-at-scale:7}
\emph{Programmation.} Vérifier \texttt{rolling\_sum\_kernel(50)} contre une version par somme cumulée sur le tableau entier, pour 100\,000 valeurs aléatoires,
et des tailles de bloc 1, 7, 50, 1\,000 et la longueur totale. Quel est le plus grand écart~?
\end{exercise}

\begin{exercise}[$\star\star\star$]\label{exo:pl:python-at-scale:8}
\emph{Trouver la faille.} «~Notre tâche de variables est lente, alors nous avons découpé la journée en huit parties, avec un pool de huit threads, chacun exécutant
la boucle Python de calcul des variables sur sa partie, puis nous avons concaténé les résultats.~»
\end{exercise}

\section{Problème~: trois minutes ou trois secondes}

\begin{problem}[{Problème du week-end --- un signal sur une année d'événements}]\label{pb:pl:python-at-scale:1}
Le signal du chapitre (une moyenne pondérée exponentiellement du volume signé), ses quatre implémentations et les mesures de
\texttt{bench\_pyscale.py}.

\medskip\noindent\textbf{Partie I --- L'interpréteur.}
\begin{enumerate}
\item Que fait l'interpréteur par événement dans la boucle sur les enregistrements qu'il ne fait pas dans l'expression sur tableaux~?
\item Combien de nanosecondes par événement coûte chacune des quatre méthodes~?
\item Quelle est l'accélération de l'expression sur tableaux par rapport à chaque boucle~?
\item Pourquoi les réponses de la boucle et du filtre sont-elles identiques au dernier bit~?
\item Quels types de boucle ne peuvent pas s'écrire comme des opérations sur tableaux, et où vont-elles~?
\end{enumerate}

\medskip\noindent\textbf{Partie II --- Mémoire.}
\begin{enumerate}[resume]
\item Combien d'octets une ligne coûte-t-elle pour une colonne de symboles en objets, en catégorie, en chaînes Arrow et en unicode à largeur fixe~?
\item Pourquoi la catégorie est-elle la plus petite~?
\item Que coûte le calcul du signal sur le tableau entier en mémoire par événement, et d'où provient ce coût~?
\item Combien d'événements tient un budget de 1.5~Go d'un coup~?
\item Qu'est-ce que le \omterm{def:pl:python-at-scale:chunked}{pic de mémoire résidente}, et pourquoi c'est lui, et non la moyenne, qui décide si un travail tient~?
\end{enumerate}

\medskip\noindent\textbf{Partie III --- Blocs et cœurs.}
\begin{enumerate}[resume]
\item Pourquoi le résultat par blocs est-il égal au résultat complet, et quel état est transporté~?
\item Quelle taille de bloc est la plus rapide ici, et quel est son pic de mémoire~?
\item Pourquoi les blocs de mille événements sont-ils plus lents~?
\item Que gagnent deux threads sur une boucle Python et sur un tri numpy~?
\item Que coûte le transfert de 64~Mo vers un autre processus par sérialisation pickle, et par mémoire partagée~?
\end{enumerate}

\medskip\noindent\textbf{Partie IV --- Le verdict.}
\begin{enumerate}[resume]
\item Énoncer le \emph{résultat nommé}~: l'accélération de l'expression sur tableaux par rapport aux boucles, la taille de bloc qui minimise le temps et
son pic de mémoire, et le plus grand historique qui tient entier dans le budget.
\item Combien de temps prendrait une année de 250 jours de deux cents millions d'événements par chaque méthode, sur un seul cœur~?
\item Comment utiliser huit cœurs pour cette année sans payer le \omterm{def:pl:python-at-scale:processes}{coût de sérialisation}~?
\item Parmi les leçons du chapitre, lesquelles un interpréteur sans verrou global modifierait-il, et lesquelles non~?
\item En une phrase~: qu'est-ce qui rend Python assez rapide pour les \omterm{def:pl:capturing-and-storing-tick-data:tick}{données tick}~?
\end{enumerate}
\end{problem}

\section{Questions d'entretien}

\begin{interviewq}[$\star$ \iqroles{chercheur, développeur}]\label{iq:pl:python-at-scale:1}
Pourquoi une boucle Python sur un tableau numpy est-elle lente, et que faites-vous à la place~?
\end{interviewq}

\begin{interviewq}[$\star\star$ \iqroles{chercheur, développeur}]\label{iq:pl:python-at-scale:2}
Comment calculeriez-vous une moyenne mobile exponentielle sur un milliard de valeurs en Python~?
\end{interviewq}

\begin{interviewq}[$\star\star$ \iqroles{développeur}]\label{iq:pl:python-at-scale:3}
Qu'est-ce que le \omterm{def:pl:python-at-scale:gil}{verrou global de l'interpréteur} (GIL), et quand n'a-t-il pas d'importance~?
\end{interviewq}

\begin{interviewq}[$\star\star$ \iqroles{chercheur}]\label{iq:pl:python-at-scale:4}
Votre dataframe utilise 40~Go de mémoire. Où regardez-vous en premier~?
\end{interviewq}

\begin{interviewq}[$\star\star\star$ \iqroles{développeur, chercheur}]\label{iq:pl:python-at-scale:5}
Concevez un travail qui calcule des variables sur cinq ans de \omterm{def:pl:capturing-and-storing-tick-data:tick}{données tick} sur une machine de 32~cœurs avec 128~Go de mémoire.
\end{interviewq}
